\documentclass[11pt]{article}
\usepackage{icmlko}
\begin{document}

\title{여섯 번 재서 여섯 번 다 $0$ --- 자를 바꿔 다시 열어도 없는 것은
없다}
\author{월드모델 실험실}
\date{노트 455 · 456 · 2026년 8월 2일}
\maketitle

\begin{abstract}
두 갈래를 닫는다. 둘 다 \textbf{그럴듯한 기계}가 있었고 둘 다
\textbf{정확히 아무 일도 안 일어났다}.

\textbf{하나 --- 갇힌 도메인의 행은 남을 해치나.} 노트 454 가 시장팝업
$249$행이 $18$칸뿐임을 보였다. 학습에 든 $123$행은 서로 거의 겹치니
나무에 \emph{가짜 순도 높은 갈래}를 만들어 줄 법했다. 학습에서만 빼고
채점은 그대로 두었다.

\begin{center}
\begin{tabular}{lrr}
\hline
 & 판 차 & 날짜 통제 차 \\
\hline
$-$ 시장팝업 & $\mathbf{+0.0000}$ & $+0.0002$ \\
$-$ 팝업 & $-0.0017$ & $-0.0031$ \\
$-$ 둘 다 & $-0.0014$ & $-0.0031$ \\
\hline
\end{tabular}
\end{center}

\textbf{원핫이 이미 그 도메인을 따로 떼어 놓고 있었다.} 갇힌 행은 갇힌
채로 자기 가지에 앉아 있고 나머지 열 도메인에 안 닿는다.

\textbf{둘 --- K$=$64 를 자를 바꿔 다시 연다.} 노트 431 이 K$=$64 를
기각한 근거는 \textbf{LODO $3/11$} 이었다. 그런데 노트 441 이
\emph{LODO 는 집 밖 효과를 원리상 못 본다}는 것을 보였고, K$=$64 의
\emph{원래} 근거는 집 밖 짝 하나(비게임 앱 되뽑기 $1.000$)였다 ---
\textbf{기각에 쓴 자가 그 종류의 효과를 못 재는 자였다.} 이제 짝이 둘이니
제대로 다시 잰다.

\begin{center}
\begin{tabular}{lrrrr}
\hline
 & 판 & 날짜 통제 & KR 만화 & 비게임 앱 \\
\hline
K$=$32 (챔피언) & $+0.4870$ & $+0.5274$ & $+0.5503$ & $+0.2614$ \\
K$=$64 & $+0.4872$ & $+0.5270$ & $+0.5512$ & $+0.2607$ \\
\hline
차 & $+0.0002$ & $-0.0004$ & $+0.0009$ & $-0.0007$ \\
\hline
\end{tabular}
\end{center}

\textbf{네 계기 전부 잡음 안이다.} 노트 431 의 판정은 그대로 서는데,
\emph{더 나은 이유로} 선다.
\end{abstract}

\begin{figure}[t]
\centering
\includegraphics[width=.92\textwidth]{figs/f1.pdf}
\caption{여섯 측정이 전부 잡음 띠 안에 있다. 점선 둘은 이번 사이클에
실제로 채택된 것들($+0.0041$ 계절 · $+0.0046$ 특기 칸)이라 크기 비교용
이다.}
\end{figure}

\section{내 기계 추측이 두 번 다 틀렸다}

\textbf{시장팝업.} ``거의 같은 행 $123$개가 나무를 오염시킨다''는 그림은
틀렸다. 원핫이 그 도메인에 자기 열을 주므로 나무는 그 $123$행을 한
가지에 몰아넣고 끝낸다 --- \emph{오염될 통로가 없다}. 덤으로 하나 더
확인됐다: 시장팝업은 학습에서 빠져도 자기 점수가 $-0.0055$ 밖에 안
떨어진다($+0.1499 \to +0.1443$). 노트 445 의 \emph{``가진 게 전부 빌린
것''}이 세 번째로 같은 값을 낸다.

\textbf{K$=$64.} 자루를 두 배로 늘리면 예보 잡음이 줄어 집 밖에서 이득이
날 법했다. 안 난다. K$=$32 에서 이미 포화다.

\section{닫는 법이 달라졌다}

이 노트의 값은 두 $0$ 자체가 아니라 \textbf{닫는 방식}에 있다.

노트 431 은 K$=$64 를 \emph{한 짝의 되뽑기 $1.000$} 으로 열었다가
\emph{LODO $3/11$} 로 닫았다. 둘 다 지금 기준으로는 못 쓰는 근거다 ---
앞은 짝이 하나뿐이라 못 쓰고(노트 441), 뒤는 자가 그 효과를 못 봐서
못 쓴다(노트 441). 그런데도 \textbf{결론은 맞았다.} 우연이었다.

지금은 근거가 셋이다 --- 판 · 날짜 통제 판(노트 446) · 집 밖 두 짝
(노트 441 · 444). \textbf{같은 결론을 우연이 아니게 만드는 데 사이클
하나가 들었다.}

\begin{quote}
\textbf{조항} --- 자(계기)가 바뀌면 그 자로 닫았던 갈래는 \emph{다시 열
수 있다}. 다시 열어 같은 답이 나오면 그때 비로소 확정으로 적는다.
K$=$32 를 여기서 \textbf{확정}으로 적는다.
\end{quote}

\section*{남은 것}

\begin{itemize}
\item \textbf{같은 잣대로 다시 열 것들.} 노트 431 이 LODO 로 닫은 것이
K$=$64 말고 또 있나 --- 노트 432 의 깊이 $12$ 와 lr $0.10$ 이 같은
자로 되돌려졌다. 깊이 $12$ 는 집 밖 짝 둘로 다시 볼 만하다.
\item \textbf{특기 칸 문턱 $\le 5$} --- 선별에서 $+0.0018$ 로 문턱 아래라
접어 뒀다. 관문을 돌리면 정해진다.
\item \textbf{시장팝업 태거} --- 모형 쪽은 이제 할 게 없다(노트 454 ·
455). 사람 태깅 스무 건에서 확인할 것은 \emph{몇 칸으로 갈리는가}다.
\end{itemize}

\end{document}
